\section{Security}\label{sec:offline:proof}

\subsection{Security analysis}

Let $\prot$ denote the protocol from Section~\ref{sec:construction}. We prove
that there exists a simulator $\simu$ such that, for every adversary $\adv$,
the joint execution of $\prot$ with $\adv$ is computationally
indistinguishable from the interaction of $\adv$ with ideal functionality
$\F$ via $\simu$. We proceed via a sequence of hybrid games, each
computationally indistinguishable from the last. Throughout, we call a user
\emph{corrupt} if both their device and secure element are compromised; the
more general split-corruption setting follows analogously. We also simplify
exposition by treating the randomness $\rho$ produced by $\F$ during
$\fwithdraw$ and $\fpay$ as the output of $\com(\ck, \cdot; \cdot)$.

We rely on the following security properties of the underlying primitives:

\begin{enumerate}

  \item[(P1)] The commitment scheme is binding: no PPT adversary can find
    two openings of the same commitment to distinct messages.

  \item[(P2)] The encryption scheme is IND-CPA secure: ciphertexts are
    computationally indistinguishable from random ones for any PPT
    adversary.

  \item[(P3)] The blind signature scheme used by $\cb$ is EUF-CMA secure: no
    PPT adversary can produce a valid signature without querying the
    signer.

  \item[(P4)] The credential signature scheme used by $\reg$ is EUF-CMA
    secure.

  \item[(P5)] The zero-knowledge proofs and signatures of knowledge are
    simulation-extractable in the random oracle model: a valid witness can
    be extracted from any accepted proof, and simulated proofs are
    indistinguishable from real ones.

\end{enumerate}

\paragraph{Game 0.} The real-world execution of $\prot$.

\paragraph{Game 1.} $\simu$ runs $\prot$ internally on behalf of all honest
parties, forwarding inputs and distributing outputs. This is identically
distributed to Game 0.

\paragraph{Game 2.} $\simu$ introduces a dummy $\F$ as a communication relay
between honest parties, who now communicate via $\F$ rather than directly.
This is identically distributed to Game 1.

\paragraph{Game 3.} $\F$ handles $\fcreate$ queries from honest users. When
an honest user $\usr$ initialises an account of value $v$ with intermediary
$\Int$, $\simu$ calls $\F$ with $(\fcreate, \usr, v)$ in place of running
$\prot$. The two are identical: $\prot$ always accepts when the user has no
prior account, and $\F$ performs the same check.

\paragraph{Game 4.} As Game 3, extended to corrupt users initiating account
creation. $\simu$ calls $\F$ and mirrors any abort to $\Int$. $\F$ aborts
exactly when $\prot$ does, on attempts to register duplicate accounts, so no
new distinguishing event arises.

\paragraph{Game 5.} $\F$ handles $\fregister$ queries. $\simu$ calls
$(\fregister, \usr, \se, \dev)$ whenever $\reg$ receives a valid
registration request. $\F$ rejects only duplicate registrations, matching
$\prot$.

\paragraph{Game 6.} $\F$ handles $\fwithdraw$ queries from honest users.
When an honest $\usr$ withdraws value $v$, $\simu$ calls $(\fwithdraw,
\Int, v)$ and learns $(\usr, v)$. $\simu$ produces a simulated withdrawal
transcript in which $\comm$ is a uniformly random commitment, $\cipher_0
\leftarrow \enc(\audpk, \uid; \rho_0)$ is a well-formed encryption of
$\uid$, $\cipher_2$ is a random ciphertext, and $\Pi$ is a simulated proof.
By (P1), (P2), and (P5), this transcript is computationally
indistinguishable from one produced by $\prot$.

\paragraph{Game 7.} $\F$ handles $\fwithdraw$ queries from corrupt users.
When $\usr$ submits a withdrawal, $\simu$ extracts the randomness $s$ from
the attached proof using (P5), computes $\rho \leftarrow \com(\ck, (\uid,
v); s)$, and calls $(\fwithdraw, \Int, v)$ with $\rho$. Both $\F$ and
$\prot$ reject if and only if $\comm \leftarrow \com(\ck, (\uid, v); s)$ is
a duplicate, which would imply a duplicate $\rho$, so the accounting checks
agree.

\paragraph{Game 8.} $\F$ handles $\fpay$ queries from honest payers. We
first deal with the first payment after withdrawal, the case $i = 0$, then
generalise.

\emph{Case $i = 0$.} $\usr$ pays $\usr'$ by calling $(\fpay, \rho, v,
\usr')$.

\begin{itemize}

  \item If $\usr'$ is honest: $\simu$ learns $(\fpay, v, \rho, \rho')$ from
    $\F$, and produces a simulated payment $\py_0 = (v, \comm_0, \comm_1,
    \cipher_0, \sigma_0, \Gamma_0, \Sigma)$ with random commitments
    $\comm_b \leftarrow \com(\ck, r_b; s_b)$ for fresh $r_b, s_b$, a random
    ciphertext $\cipher_0$, and simulated proofs $(\sigma_0, \Gamma_0,
    \Sigma)$, possible by simulation-extractable security of the NIZK
    system (P5).

  \item If $\usr'$ is corrupt: $\simu$ additionally learns the identities
    of both parties. It computes $\rho' \leftarrow \com(\ck, \uid'; s')$
    from the commitment $\comm_1$ that $\usr'$ supplies in $\prot$, calls
    $\F$, and sets $\cipher_0 \leftarrow \enc(\audpk, \uid; \rho_0)$.

\end{itemize}

In both cases $\simu$ records $\mathbb{R}[(\comm_0, \comm_1)] \gets (\rho,
\rho')$ and $\mathbb{C}[(\rho, \rho')] \gets (\comm_0, \comm_1)$.

\emph{Case $i > 0$.} $\simu$ recovers $(\comm_0, \ldots, \comm_i)$ from
$\mathbb{C}[\vec{\rho}]$, produces a simulated payment $\py_i$ with
simulated proofs and random ciphertexts (or $\cipher_i \leftarrow
\enc(\audpk, \uid_i; \rho_i)$ if $\usr'$ is corrupt), and updates the maps
as needed. By IND-CPA security of the encryption scheme (P2) and
simulation-extractable security of the NIZK scheme (P5), simulated payments
are indistinguishable from real ones.

\paragraph{Game 9.} $\F$ also processes $\fpay$ queries from corrupt
payers. We show that if $\F$ rejects a payment, an honest $\usr'$ in
$\prot$ also rejects it, again showing the case $i = 0$ before generalising.

Let $\py_0 = (v, \comm_1, \comm_0, \cipher_0, \sigma_0, \Gamma_0, \hist_0)$
be the payment produced by corrupt $\usr$. $\simu$ computes $\rho \leftarrow
\com(\ck, (\uid, v); s)$ from $\comm_0$ and calls $(\fpay, \rho, v,
\usr')$.

\emph{$\usr'$ honest.} If $\usr'$ accepts $\py_0$ but $\F$ rejects, exactly
one of the following holds:

\begin{enumerate}[label=\alph*)]

  \item $\tokens[\rho] = \bot$: no prior $\fwithdraw$ matches $\rho$. By
    (P5), $\simu$ extracts a blind signature $\psi$ and message triple
    $(\uid, v, s)$ from $\Sigma$. Since $\cb$ never issued $\psi$, no
    corresponding withdrawal exists, so outputting $(\psi, \uid, v, s)$
    breaks EUF-CMA of the blind signature scheme (P3).

  \item $\tokens[\rho] = \langle \usr, v^* \rangle$ with $v^* \neq v$: we
    can extract $(\ck, \uid', s', \comm_0)$ from $\Sigma$ such that
    $\opn(\ck, \uid'; s', \comm_0) = 1$, giving $\rho \leftarrow \com(\ck,
    \uid'; s')$ with value $v$. Combined with $\rho \leftarrow \com(\ck,
    \uid^*; s^*)$ with value $v^*$ from the prior withdrawal, $v \neq v^*$
    yields two distinct openings of $\rho$, contradicting binding of the
    commitment scheme (P1).

  \item $\tokens[\rho] = \langle \usr^*, v \rangle$ with $\usr^* \neq
    \usr$: if $\usr^*$ is corrupt, the same two representations of $\rho$
    as in (b) give $\uid \neq \uid^*$, again contradicting (P1). If
    $\usr^*$ is honest, $\comm_0 \leftarrow \com(\ck, \uid^*; s^*)$ was
    chosen by $\simu$, so $\usr$ cannot have guessed it except with
    negligible probability; if $\usr'$ accepts $\py_0$, then $\sigma_0$ is
    a valid signature of knowledge proving possession of a credential on
    $(\uid, \sk)$, and by (P5) $\simu$ extracts credential $\varphi$ and
    pair $(\uid, \sk)$. Since $\uid$ was assigned randomly by $\simu$ and
    was not registered, this gives a forgery of the credential scheme
    (P4).

  \item $\usr$ is not registered: if $\usr'$ accepts, $\sigma_0$ verifies,
    and by (P5) $\simu$ extracts a credential forgery $(\varphi, (\uid,
    \sk))$, contradicting (P4).

\end{enumerate}

\emph{$\usr'$ corrupt.} $\simu$ sends $\accept$ to $\F$ on $\usr'$'s
behalf, provides $\rho' \leftarrow \com(\ck, (\uid', v); s')$ from
$\comm_1$, and extends the maps.

\emph{$i > 0$.} If $\mathbb{R}[(\comm_0, \ldots, \comm_i)]$ is not yet
recorded, $\simu$ fills in the missing intermediate payments by extracting
participant identifiers from the proof $\Gamma_j$ at the first unrecorded
index $j$ using (P5), calling $\fpay$ for each, and extending the maps until
$j = i - 1$. The cases above still apply, replacing $\py_0$ with $\py_i$ and
$\rho$ with $(\rho_0, \ldots, \rho_i)$. Since the probability that an honest
$\usr'$ accepts a payment that $\F$ rejects is negligible under (P1)-(P5),
this game is indistinguishable from Game 8.

\paragraph{Game 10.} $\F$ handles $\fdeposit$ queries from honest users.
$\simu$ recovers $(\comm_0, \ldots, \comm_n)$ from $\mathbb{C}[\vec{\rho}]$
and produces a simulated deposit $(\tau_n, \sigma_n, \hist_n)$ with
simulated proofs and random ciphertexts. Since $\F$ rejects a deposit only
when the token does not belong to $\usr$ or has been previously deposited,
and an honest user never deposits a token they do not own, the only case
requiring argument is double-spend detection.

Each $\rho_i$ is determined by $\comm_i$ via $\rho_i \leftarrow \com(\ck,
(\uid_i, v); s_i)$, so the mapping $\rho_i \leftrightarrow \comm_i$ is
one-to-one by binding of the commitment scheme (P1). $\rho_0$ is fixed by
$\comm_0$, which is fixed by the blind signature issued by $\cb$, which
binds $(\uid_0, v, s_0)$. A corrupt user cannot alter $\comm_0$, so any two
deposits sharing $\vec{\rho}[0]$ are double spends and will be detected by
$\F$.

\paragraph{Game 11.} $\F$ now also handles $\fdeposit$ queries from
corrupt users. When $\tokens[\vec{\rho}] \neq \langle \usr, v \rangle$,
exactly one of the following holds, each handled by the same argument as in
Game 9 with the payment in question replaced by the deposit in question.

\begin{enumerate}[label=\alph*)]

  \item $\tokens[\vec{\rho}] = \bot$: extracting $\Sigma$ via
    simulation-extractability of the NIZK scheme (P5) gives a blind
    signature forgery, contradicting unforgeability of the signature
    scheme.

  \item $\tokens[\vec{\rho}] = \langle \usr^*, v^* \rangle$ with $v^* \neq
    v$: two valid openings of $\rho_0$ under $\opn$ with distinct values
    violate binding of the commitment scheme (P1).

  \item $\tokens[\vec{\rho}] = \langle \usr^*, v \rangle$ with $\usr^*
    \neq \usr$: either binding of the commitment scheme (P1) is violated,
    or extracting $\sigma_0$ via (P5) yields a credential forgery,
    contradicting EUF-CMA of the credential signature scheme (P4).

  \item $\usr$ unregistered: $\sigma_0$ verifying implies a credential
    forgery, contradicting unforgeability of the credential signature
    scheme.

\end{enumerate}

If the token has been transferred multiple times before deposit, $\simu$
fills in unrecorded payments using $\Gamma_j$ as in Game 9, and the same
cases apply inductively. If another deposit chain sharing
$\vec{\rho}[0]$ has already been recorded, the same argument as in Game 10
shows the initial commitment cannot be altered.

\paragraph{Game 12.} $\F$ handles $\faudit$ queries. The auditor calls
$(\faudit, \rho)$ and expects the set of double spenders $\cDS$. We show
that the outputs of $\F$ and $\prot$ agree except with negligible
probability.

In $\prot$, the double spender at index $j$ is identified by decrypting
$\cipher_j$. The commitment $\comm_j$ is fixed by Game 11's argument
(binding of the commitment scheme and EUF-CMA of the blind signature
scheme) through the chain via simulation-extractability of $\sigma_{j-1}$.
The ciphertext $\cipher_j \leftarrow \enc(\audpk, \uid_j; \rho_j)$ encrypts
the same value as is committed in $\comm_j$, as certified by proof
$\Gamma_j$; were this false, it would either yield two valid openings of
$\comm_j$ under $\opn$, contradicting binding (P1), or a proof $\Gamma_j$
with no valid witness, contradicting simulation-extractability (P5). An
honest recipient cannot be framed, as $\comm_j$ is chosen by the recipient
before being signed by the payer in $\sigma_{j-1}$; hence $\comm_j =
\comm_j'$ if and only if the owner of $\comm_j$ is the double spender.

If $\cipher_j$ decrypts to an unregistered user, then $\sigma_j$ is a valid
signature of knowledge over a credential for that identifier, contradicting
EUF-CMA of the credential scheme (P4) via simulation-extractability of the
NIZK (P5). All distinguishing events occur with negligible probability, so
Game 12 is indistinguishable from Game 11. This completes the proof.

\subsection{Discussion}\label{sec:securityDiscussion}

The above proves security in the cryptographic sense, under the standard
assumption that a real-world deployment realises each ideal primitive it
relies on. We close by discussing, informally, some further considerations
that arise when deploying the system in practice.

\paragraph{Revocation.} A practical deployment requires the ability to
revoke users. This is complicated by two factors: first, a user's privacy
should not be violated by revocation itself, meaning that no new
information about their past transactions should become available once they
are revoked; second, and more fundamentally, revocation interacts poorly
with the offline setting's inherent lack of up-to-date information. Consider
a payee who goes offline just before a payer is revoked: the payee's
revocation metadata will not reflect the payer's revocation, so the payee
will accept the revoked payer's token, discover the problem only later, and
may unwittingly propagate it further by spending the token onward.

One effective mitigation is a time-to-live (TTL) for offline tokens,
enforced either by capping the number of spending hops a token may
undergo or by bounding the time elapsed since the original withdrawal.
A malicious spender cannot be trusted to enforce such a policy, but payees
are well placed to do so themselves, by refusing payments whose TTL does not
meet policy; the policy can be tuned to the value of the transaction, with
higher-value transactions given a shorter TTL. Capping the value of
individual offline transactions is a complementary mitigation, and combined
with a TTL can make double spending financially unattractive even before
any detection takes place. We leave a fuller treatment of revocation and
its interaction with TTL policies to future work.

\paragraph{Custody of the secret key.} Perhaps the most important
requirement for a real-world deployment is that the user not be in direct
possession of secret key $\usk$. If a user does hold $\usk$ directly, they
can generate an arbitrary number of seemingly valid payments spending the
same token, each of which will pass every check available to its offline
recipient; the impossibility of detecting double spending before
reconciliation is inherent to the offline setting, not a flaw in a
particular recipient's checks. Payments remain auditable, since each
carries a ciphertext that provably encrypts the payer's identity, and this
holds even under \emph{cascade} double spending, where a payee further
spends an already double-spent token: every perpetrator in the chain is
still identified. This threat of identification is a meaningful deterrent
only if detection can be attributed to malice rather than to key
compromise, since a user whose device is hacked and a user who
deliberately double spends leave the same digital trail; in this scenario,
cryptography can support double-spending resistance but cannot enforce it,
and the remaining gap must be closed at the governance level, for instance
through insurance or liability for compromised devices, with TTL
enforcement again limiting the exposure.

Mediating access to the secret key through a trusted element that only ever
produces a single payment per token closes off this class of attack, at the
cost of introducing others: an adversary who can clone or roll back the
secure element can still double spend, though now at a real hardware cost
rather than none, and with the same reporting incentives making
undetected repeat use harder to sustain. A properly implemented system also
needs a recovery mechanism for lost or destroyed devices, which itself
creates an avenue for abuse if a user falsely claims their device is lost in
order to retain and later use it. Introducing a waiting period before a
replacement device is issued, timed so that it exceeds the TTL of the
locked tokens, together with policies that flag frequent recoveries, limits
this risk without giving up the usability benefits of recovery.

\paragraph{Instantiating the secure element.} Modern mobile platforms
provide trusted execution environments that protect app secrets from a
malicious device owner, but app developers cannot directly access
TrustZone~\cite{trustzone} or Secure Enclave~\cite{secureenclave}: these
hardware features are exposed only through higher-level key-management
APIs. While the secret key itself remains hidden even from a malicious app,
an attacker who achieves arbitrary code execution, whether by exploiting the
app or by rooting the device, can still query the key manager to spend a
token repeatedly, since the confidentiality of the key does not guarantee
the integrity of the surrounding application logic. Smart cards offer a
stronger guarantee: platforms such as Javacard~\cite{javacard} or
MIFARE~\cite{mifare} provide tamper-resistant storage together with a
dedicated, isolated execution environment, so that both the secret key and
the logic that uses it are protected. No hardware security is unconditional,
and an attacker willing to invest sufficient resources may eventually
extract a card's embedded secret, reverting to the threat model discussed
above; deploying the system is therefore ultimately a question of
governance, weighing the cost of rolling out smart cards, the total offline
value permitted under system-wide restrictions such as TTL, and the
expected cost and payoff of attacking a card.
